problem:  
Let \((M^{1+n}, g)\) be a globally hyperbolic Lorentzian manifold that is **asymptotically flat**, in the sense that outside a compact region there exist coordinates on each Cauchy hypersurface \(\Sigma_t \simeq \mathbb{R}^n\) with \(g = \eta + h\) and \(h = O(|x|^{-1-\epsilon})\) for some \(\epsilon > 0\), where \(\eta\) is the Minkowski metric.  
Let \((N,h)\) be a compact Riemannian manifold of nonpositive curvature, and consider the finite-energy wave map
\[
\square_g \phi + A(\phi)(d\phi,d\phi) = 0, \qquad \phi : M \to N.
\]

In the flat case \(g=\eta\), the **Bahouri–Gérard profile decomposition** provides the analytic foundation for the conjecture that solutions asymptotically decompose into finitely many harmonic-map solitons and dispersive radiation (the soliton resolution conjecture).  

**Problem (Stability of Soliton Resolution under Geometric Perturbations):**  
Assume that for wave maps on Minkowski space \((\mathbb{R}^{1+n},\eta)\), the soliton resolution conjecture holds in the energy-critical dimension \(n=2\).  
Is the soliton resolution mechanism *stable* under small asymptotically flat Lorentzian perturbations of the metric?  

More precisely, does there exist \(\delta > 0\) such that if \(\|g - \eta\|_{C^3} < \delta\) in weighted Sobolev norms, then every finite-energy wave map \(\phi:(M,g)\to N\) admits, as \(t\to\infty\), an intrinsic decomposition
\[
E[\phi] = E[\phi_{\mathrm{rad}}] + \sum_{j=1}^J E[\psi^{(j)}] + o(1),
\]
where each \(\psi^{(j)}\) is a Lorentz transform of a harmonic map (a geometric soliton bubble), and the radiative component \(\phi_{\mathrm{rad}}\) scatters to a solution of the linearized wave map equation around a constant map with respect to \(g\)?  

Equivalently: is the soliton resolution structure for wave maps **geometrically stable** under perturbations of the domain metric?

Why is it a “good’’ problem:  
This formulation isolates the genuinely new and profound aspect: the **stability** of the soliton resolution phenomenon under small geometric perturbations of the Lorentzian metric. While the soliton resolution conjecture itself is a central open problem, its persistence under curvature perturbations is unexplored—and would reveal whether the conjecture is an artifact of flat geometry or an intrinsic property of nonlinear geometric wave dynamics.  

It is precise, conceptually simple, and bridges two frontiers: nonlinear dispersive analysis and Lorentzian geometry. A positive outcome would show that asymptotic integrability of nonlinear geometric fields is not tied to special symmetries of flat space, but rather a robust geometric feature—thus linking the analytic structure of wave equations with the stability of spacetime geometry. This makes it a deeply insightful and field-connecting research problem.